\section{Foundation}

This text considers the \textit{Physical Universe} strictly as a
\textit{Conceptual Experience} of a \textit{Cognitive Subject}. Physics is a
discipline of models, and models are conceptual architectures that exist in
minds. Physical time, space, matter, and energy are not the ground reality
itself; they are traces that an un-conceptual reality leaves upon the mind of
an observer. As traces, they are not real in the way one's own existence and its
experience are real. An inquiry into the origin of reality cannot commence with physical
concepts, as entities that are not real cannot produce something real.
Concepts in general, however, are equally incapable of generating existence.
What remains stable through all existence is the behavior of logic and
mathematics. This investigation therefore traces cosmogony by means of
\textit{logical necessities}---structural requirements that must be fulfilled
for an objective Universe to be perceived as the Conceptual Experience of a
Cognitive Subject.

\subsection{Nothingness}

Let \textit{Nothingness} be defined as the absolute absence of
anything---including the absence of any capacity to \textit{instantiate}, that
is, the primordial capacity to bring an entity into existence. This absolute
absence necessarily includes the absence of physical laws, mathematical
structures, quantum fields, probability distributions, and any form of latent
potentiality\footnote{A quantum vacuum governed by field equations is not
Nothingness---it is a richly structured physical state presupposing an entire
framework of laws, which are themselves conceptual architectures requiring
precisely the cognitive foundation this inquiry seeks to establish.}.
Nothingness, thus defined, is self-sealing: the transition from non-existence
to existence requires a capacity that Nothingness, by definition, lacks.
The existence of something denies Nothingness as pre-existing everything.

\subsection{The Transconceptual}

Whatever this positively existing ground is, we can bracket it by strict
deduction from the constraints already established.

It is \textbf{not Nothingness}---this was demonstrated above. It is
\textbf{non-conceptual}---because concepts require \textit{Differentiation},
and Differentiation in turn requires concepts, namely attributes by means of
which differences can be established. Neither can exist without the other; at
the absolute origin, neither exists. It is a \textbf{Unity}---because
\textit{Plurality} requires Differentiation, which is unavailable. It is
\textbf{atemporal}---because time requires differentiable moments, and
Differentiation is unavailable.

\begin{formalism}[The Transconceptual]
    The \textbf{Transconceptual} ($\mathcal{T}$) is the positively real Ground
    of existence. It is:
    \begin{enumerate}[label=(\roman*)]
        \item Not Nothingness.
        \item Non-conceptual.
        \item A Unity.
        \item Atemporal.
    \end{enumerate}
\end{formalism}

This formalism, expressed in conceptual language, cannot grasp the
Transconceptual itself. It defines it strictly by what it is not in conceptual
terms---an apophatic determination that marks the boundary of what formal
reasoning can reach.

\subsection{Receptivity and Information Substance}

Before any concept, any experience, any perception can arise, there must be
something that can potentially be experienced. Let this be called
\textit{Information Substance}: a raw, pre-conceptual form of truth that exists
as a pure \textit{potential to inform}.

\begin{formalism}[Information Substance]
    \textbf{Information Substance} is the manifestation of truth from the
    Transconceptual as a potential to inform. It is information necessarily
    bound to truth.
\end{formalism}

Falsehood requires Differentiation and comparison, neither of which is
available at the origin. Truth is not one mode among alternatives; it is the
only mode of existence.

Information can only exist if it can inform. Information that cannot inform is
not information---it is unable to exist as such. Therefore, Information
Substance can only exist alongside a corresponding \textit{Lack of Knowledge}.
Moreover, a Lack of Knowledge must exist in a logical step prior to Information
Substance, for the latter to be able to exist at all. A more precise term for
this Lack of Knowledge is \textbf{Receptivity}. An entity that is receptive for
information must be void of that particular information. It is therefore the
Transconceptual, by its nature beyond concepts, that generates the Receptivity
for Information Substance.

\begin{formalism}[Receptivity]
    \textbf{Receptivity} is the structural precondition for Information
    Substance to exist. It is generated by the Transconceptual's non-conceptual
    nature: an entity beyond concepts is inherently receptive for conceptual
    determination.
\end{formalism}

The Receptivity, however, is restricted by the Information Substance which it
has already received. This restriction separates it from the Transconceptual,
which is void of restrictions. The Receptivity is therefore a \textit{separate
entity}. Since it is a consequence of the Transconceptual, this separation
constitutes the first instance of \textit{Activity} and \textit{Passivity}: the
Transconceptual acts; the Receptivity receives. With each further Information
Substance received, the Receptivity undergoes change: its openness is
progressively shaped by the truths it has absorbed, while remaining receptive
for truths it has not yet received.

The primary truth, the first Information Substance, may only carry information
about the Transconceptual and the Receptivity, because nothing else exists.

\begin{formalism}[The Primordial Discernible $\delta_\bullet$]
    Let ``I'' be a logical placeholder for the Identity of the absolute Ground
    Truth---the Transconceptual---in its complete isolation. Let ``You'' be a
    logical placeholder for the Identity of the primordial Receptivity. Then,
    the first informational substance to exist is $\delta_\bullet$, namely:

    \begin{equation}
         \delta_\bullet := \text{``I am active, You are passive''}
    \end{equation}
\end{formalism}

The bracket established that concepts and Differentiation are mutually
dependent: neither can exist without the other, and at the origin neither
exists. The primordial discernible $\delta_\bullet$ breaks this co-dependency.
The distinction between Active and Passive---between the Transconceptual and
its Receptivity---is the first \textit{Differentiation}. It arises not through
conceptual reasoning (which would require Differentiation already in place) but
through the Transconceptual's generative act, which is not bound by the
conceptual constraints it precedes.

We can derive \textit{that} the co-dependency was broken: Differentiation
exists, yet it could not have arisen from within the conceptual domain, which
presupposes it. We can derive \textit{by what}: the only entity beyond the
conceptual domain is the Transconceptual. What we cannot derive is \textit{how}
---because any description of the mechanism would itself be conceptual, and the
act precedes conceptuality. This limitation is not an omission; it is a
prediction of the framework. At the boundary of the Transconceptual, the \textit{that}
and the \textit{by what} are accessible; the \textit{how} is structurally
inaccessible.

\subsection{Imperfection and Illusion}

A methodological clarification is in order. There is no concept of a
\textit{perspective} from the Transconceptual, because there are no concepts at
the Transconceptual. Any description of how things ``appear from the
Transconceptual'' is a category error. The Receptivity is separate from the
Transconceptual, and therefore not absolute Truth. Yet we cannot assign a
\textit{How} to any relation of the Transconceptual toward the Receptivity. We
cannot state that the Transconceptual, as absolute Truth, assigns Falsehood,
Imperfection, or even Separation to the Receptivity. This confronts us with
the hard boundary of our conceptual understanding: we are aware of a truth,
and we know that we cannot grasp it.

From within the Receptivity, however, the relation presents itself as the
difference between truth at a given logical step and absolute truth. The
characterization of the Receptivity as imperfect and separated is an
implication that follows directly from its not being the Transconceptual. And
yet, from the perspective of conceptuality, the Receptivity shares one
property with the Transconceptual: existence itself.

The following sections work toward \textit{Consciousness} and
\textit{Concepts}---the faculties by which Information Substance becomes
intelligible.

\subsection{Methodological Caveat}

It may appear paradoxical that we employ concepts---including the term
`Transconceptual' itself---to circumscribe entities that exist beyond
conceptualization. This paradox is inherent and unavoidable. Any statement
made herein pertains strictly to the \textit{shadow} that the process of
manifestation casts upon our conceptual understanding. The present
investigation aims to explore how far the analysis of these shadows may grant
\textit{insight} into the nature of the reality that produces them, while
acknowledging that the shadows remain shadows.

\subsection{Summary of Foundational Posits}

\begin{formalism}[Demarcation]
    The foundational posits of this investigation are:
    \begin{description}
        \item[The Transconceptual ($\mathcal{T}$):] The positively real,
        undifferentiated Ground of existence beyond all conceptual
        determination. Absolutely true.

        \item[Receptivity:] The structural precondition for Information
        Substance; generated by the Transconceptual's non-conceptual nature.
        Restricted by received Information Substance, and thereby separate
        from the unrestricted Transconceptual. Shares with the Transconceptual
        the property of existence.

        \item[Information Substance:] The manifestation of truth from the
        Transconceptual as a potential to inform. Necessarily bound to truth,
        as falsehood requires Differentiation and comparison, neither of which
        is available at the origin.
    \end{description}
\end{formalism}

Notably, the nature of the Transconceptual as existing beyond conceptuality
allows only for an apophatic definition: using concepts, we can say what it is
not, but we cannot say what it is. Nothingness, on the other hand, acts as the
lower boundary. Together with Nothingness, these three posits establish the
ontological ground on which the formal machinery of subsequent sections
operates. The derivation of Concepts and of the Cognitive Subject requires
faculties to be established in later sections.

\subsection{Demarcations}

With the foundational posits established, we distinguish the Transconceptual
from several superficially adjacent concepts. It is not Kant's
\textit{Noumenon}~\cite{Kant1781}, which is defined purely negatively and
without causal efficacy---functionally Nothingness for the purposes of
cosmogony. It is not Hegel's \textit{Absolute}, which realizes itself
\textit{through} concepts in a dialectical process unfolding in historical
time---the Transconceptual is non-conceptual and atemporal by derivation. It is
not Spinoza's \textit{Substance}, which is conceptually determinable through
attributes and modes---the Transconceptual admits no internal structure
accessible to reason. It is not an Aristotelian \textit{First Cause}, which
presupposes causal chains and temporal sequence---neither of which is available
at the origin. Nor is it a gap-filler invoked to forestall further scientific
or logical exploration; it is the very limit at which conceptual understanding
structurally terminates.

The Transconceptual, as an absolute, unrestricted, generative ground beyond
conceptualization, invites theological terminology. However, the
Transconceptual is not posited, but derived. It is derived as a logical
argument against Nothingness. The reality of Activity and Passivity between the
Transconceptual and the Receptivity is derived from the breaking of the
co-dependency of concepts and Differentiation. Its properties are bracketed
apophatically, not revealed. No article of faith is required.

It would be an impoverishment, however, to dismiss philosophical traditions
that used the precise term of a ``deity'' to express and arrive in the vicinity
of what formal reasoning reaches here: the \textit{prima causa}. The apophatic
determination of the Transconceptual, the character of absolute unity, and its
generative role stand perhaps closest to the works of Ibn Arabi
(1165--1240)~\cite{IbnArabi1203}. The convergence of the Transconceptual with
his \textit{al-Dh\=at} extends to the structural principle that self-disclosure
requires a receiving capacity distinct from the Essence itself.

The frameworks diverge at the point of \textit{wa\d{h}dat al-wuj\=ud}, the
unity of being. For Ibn Arabi, the receiving essences (\textit{al-a'y\=an
al-th\=abita}) are differentiated \textit{within} the Essence as ``realities
that have not smelled the fragrance of existence.'' This already contradicts
the nature of the Transconceptual, within which there can be no
Differentiation. Further, the derivation of the primordial discernible exposes
a decisive divergence. $\delta_\bullet$ can only assert ``I am active, You are
passive.'' Equating ``I'' and ``You'' implies Activity and Passivity as
separate within the Transconceptual, directly contradicting the pure Unity
required by its nature of existing beyond conceptuality. Ibn Arabi's
\textit{wa\d{h}dat al-wuj\=ud} cannot accommodate this.

The following sections develop a formal framework of this cosmogony. It
operates strictly on Information Substance received by the Receptivity. Its
logical operations are of a different kind than physical ones; seeking physical
explanation or convergence with physics is a category error. The framework
advances by formal necessity alone, toward the eventual derivation of
conceptual perception and, beyond it, a Cognitive Subject.
